\documentclass[11pt]{article}
\usepackage{../homework}

\profname{Dr. Stanislaw Szarek}
\classcode{MATH 423 Measure Theory}
\duedate{16 November 2026}
\assignnum{Assignment 11}

\begin{document}
\begin{problem}{2.6.55}
  Let \( E = [0,1] \times [0,1] \).
  Investigate the existence and equality of \( \int_E f \dd{m}^2 \), \( \int\limits_0^1 \int\limits_0^1 f(x,y)\dd{x} \dd{y} \), and \( \int\limits_0^1 \int\limits_0^1 f(x,y)\dd{y} \dd{x} \) for the following \( f \).
  \begin{enumerate}[label=\alph*.]
    \item \( f(x,y) = (x^2 - y^2)(x^2 + y^2)^{-2} \).
          \begin{proof}
            We first use calculus to determine the iterated integrals.
            First of all, modify the function as \( f(0,0) = 0 \), which is only a change on a null-set, so the integrals are not affected.
            Set \( x = y\tan z \), so
            \begin{equation}
            \label{eq:1}
            \begin{aligned}
              \int\limits_0^1 \int\limits_{0}^{1} \frac{x^2 - y^2}{(x^2 + y^2)^2} \dd x \dd y &= \int\limits_0^1 \int\limits_{x=0}^{x=1} \frac{y^3(\tan^2 z - 1)\sec^2z }{y^{4}\sec^4z}\dd{z}\dd{y} \\
                                                                                            &= \int\limits_0^1 \int\limits_{x=0}^{x=1} \frac{\sin^2z - \cos^2z}{y} \dd{z} \dd{y} = \int\limits_0^1 \int\limits_{x = 0}^{x = 1} \frac{-\cos 2z}{y} \dd{z} \dd{y} \\
              &= \int\limits_0^1 \frac{-1}{(1 + y^2)} \dd{y} = - \frac{\pi}{4}.
            \end{aligned}
          \end{equation}
          On the other hand, using the substitution \( y = x\tan z \), we obtain
          \begin{equation}
          \label{eq:2}
          \int\limits_0^1 \int\limits_0^1 \frac{x^2 - y^2}{(x^2 + y^2)^2}\dd{y}\dd{x} = \frac{\pi}{4}.
        \end{equation}
        This is the contrapositive of the Fubini theorem, so that we have shown \( f \not\in L^1(\mu \times \nu) \).

        It remains to show that we can actually evaluate these iterated integrals.
        Note that the positive and negative parts of \( f \) are just \( f^+ = f \text{ if } x \geq y, 0 \text{ otherwise} \), and vise-versa for \( f^- \).
        Additionally, since \( f(x,y) = -f(y,x) \), we see that the two halves of the ``triangle'' are equal,
        \begin{equation}
        \label{eq:12}
        \int f^+ \dd m^2 = \int f^- \dd m^2.
        \end{equation}
        However, since we already showed that \( f \not\in L_1 \), this means we must have both the positive and negative integrals infinite, and thus \( \int f \dd m^2 \) is undefined.
          \end{proof}
    \item \( f(x,y) = (1-xy)^{-a} \) (\( a > 0 \)).
          \begin{proof}
            Since \( f(x,y) \) is symmetric under interchange of variables, the two iterated integrals are necessarily equal.
            Moreover, the function is strictly positive for all \( (x,y) \neq (1,1) \).
            We can modify \( f(x,y) \) to take \( 0 \) at \( (1,1) \) which doesn't change the value of the integrals.
            Since this modified \( f(x,y) \) is now non-negative on all of \( E \), we can directly apply Tonelli's theorem for equality, so that all we must show is existence.

            Consider the inner integral, which can be calculated using elementary calculus.

            \begin{equation}
            \label{eq:3}
            \begin{aligned}
              I_x(y) \coloneq \int\limits_0^1 (1-xy)^{-a}\dd{x} = \begin{cases}
                \frac{(1-y)^{-a+1} - 1}{(a-1)y}, & a \neq 1 \\\\
                -\frac{\ln(1-y)}{y}, & a = 1
              \end{cases}.
            \end{aligned}
          \end{equation}
          At \( y = 0 \), the integrand is just \( 1 \), so that \( I_x(0) = 1 \).
          We split into three regions of \( a \),
          \begin{itemize}
            \item \( a \in (0,1) \): In this case, \( 1 > 1-a > 0 \), so that
                  \begin{equation}
                  \label{eq:4}
                  \begin{aligned}
                    \lim_{y\rightarrow1} \frac{(1-y)^{1-a}-1}{(a-1)y} &= \frac{-1}{a-1} \\
                    &= \frac{1}{1-a} > 0,
                  \end{aligned}
                \end{equation}
                  which is bounded from above and continuous in this region and thus integrable against \( \dd{y} \).
            \item \( a = 1 \): Here, we have
                  \begin{equation}
                  \label{eq:6}
                  -\frac{\ln(1-y)}{y},
                \end{equation}
                  which blows up exactly at \( y=1 \) but as \( y \rightarrow 1^- \) can be bounded by \( -\ln(1-y) \), which is in \( L^1(E) \), since
                  \begin{equation}
                  \label{eq:7}
                  \int\limits_0^1-\ln(1-y)\dd{y} = (1-y)\ln(1-y) - (1-y)\eval_0^1 = 1.
                \end{equation}
                  Therefore this region is also integrable.
            \item \( a > 1 \): Here, \( 1-a < 0 \), so that \( I_x(y) \) goes like
                  \begin{equation}
                  \label{eq:9}
                  \frac{-(1-y)^{1-a}}{1-a},
                \end{equation}
                  so we would need \( \int\limits_0^1 (1-y)^{1-a} \dd{y} < \infty \).
                  \begin{equation}
                  \label{eq:10}
                  \begin{aligned}
                    \int\limits_0^1 (1-y)^{1-a} \dd{y} = -\frac{(1-y)^{2-a}}{2-a} \eval_0^{1}.
                  \end{aligned}
                \end{equation}
                  If \( a < 2 \), then this clearly equals \( 0 \).
                  However, if \( a \geq 2 \), the numerator tends to \( \infty \) as \( y \rightarrow 1^- \).
                  Therefore, this region is integrable if and only if \( a < 2 \).
          \end{itemize}

          This shows that all three integrals are finite and equal for \( 0 < a < 2 \), and diverge to \( \infty \) for \( a \geq 2 \).
          \end{proof}
    \item \( f(x,y) = (x - \frac{1}{2})^{-3} \) if \( 0 < y < |x - \frac{1}{2}| \), \( f(x,y) = 0 \) otherwise.
          \begin{proof}
            Note that the region on which \( f \) is supported is two triangles, one with base \( (0,1/2 \) and the other \( (1/2,1) \), and both with height \( 1/2 \).
            It is immediately clear that the two iterated integrals should be different, since a \( y \)-slice will cut across both triangles while an \( x \)-slice will cut across only one.
            To see this, consider that the bound \( |x- 1/2| > y \) implies \( 1/2-y < x < 1/2 + y \), so we can write
            \begin{equation}
            \label{eq:13}
            \begin{aligned}
              \iint f(x,y) \dd{x}\dd{y} &= \int_{} \left( \int_{0}^{1} (x- \frac{1}{2})^{-3} \dd{x} \right) \dd{y} \\
              &= \int \left(\int\limits_0^{1/2 - y} \frac{\dd{x}}{(x-1/2)^{-3}} + \int_{1/2 + y}^1 \frac{\dd{x}}{(x-1/2)^{-3}} \right )\dd{y}\\
              &= \int \left( - \frac{1}{2y^2} - (-2) + -2 - (- \frac{1}{2y^2}) \right) \dd{y} = 0.
            \end{aligned}
          \end{equation}

          On the other hand, for \( x \neq \frac{1}{2} \) and \( x \in [0,1] \), write
          \begin{equation}
          \label{eq:14}
          \begin{aligned}
            I(x) &= \int\limits_0^{|x-1/2|}(x- \frac{1}{2})^{-3} \dd{y} \\
            &= \frac{|x- \frac{1}{2}|}{(x - \frac{1}{2})^3} = \begin{cases}
              - \frac{1}{(x - 1/2)^2}, & x < 1/2 \\
              \frac{1}{(x-1/2)^2}, & x > 1/2
            \end{cases}.
          \end{aligned}
        \end{equation}

        Then,
        \begin{equation}
        \label{eq:15}
        \iint f \dd{y} \dd{x} = \int I^+ \dd{x} - \int I^- \dd{x},
      \end{equation}
      but
      \begin{equation}
      \label{eq:16}
      \int I^+ \dd{x} = \int_{1/2}^1 \frac{1}{(x-1/2)^2} \dd{x} = \infty,
    \end{equation}
    and the same for \( I^- \), so both parts diverge and thus \( \iint f \dd{y}\dd{x} \) diverges.

    Finally, for the product measure, we can write \( \int f \dd{m^2} = \int f^+ \dd{m^2} - \int f^- \dd{m^2}\), and apply Tonelli's theorem to these non-negative parts, and appealing to the same argument as above show that \( \int f \dd{m^2} \) diverges.

    Hence, the only one of the three integrals which exists is \( \iint f \dd{x}\dd{y} = 0\), while the others diverge to infinity.
          \end{proof}
  \end{enumerate}
\end{problem}

\begin{problem}{2}
  Let \(\nu\) denote the counting measure on the set \( \mathbb{N} \) and let \( (X, \mathcal{M}, \mu) \) be a \(\sigma\)-finite measure space.
  The objective of this problem is to make as explicit as possible the product measure space \( (\mathbb{N} \times X, \mathcal{P}(\mathbb{N}) \otimes \mathcal{M}, \nu \times \mu) \eqcolon (Y, \mathcal{N}, \rho)\) and various contructs/theorems in the context of that space.
  \begin{enumerate}[label=\alph*.]
    \item Describe \( Y, A \in \mathcal{N} \) and \( \rho \) in terms of more elementary mathematical objects.
          The goal here is not to restate the definition, but to come up with an equivalent description that does not use the world ``product.''
          \begin{answer}
            \( Y \) is the space of coordinate pairs \( (n, x) \), where \( n \) is a positive integer, \( x \in X \).
            Alternatively, one can view this as a countable disjoint union of \( X \) with itself, where the first coordinate \( n \) indexes which ``copy'' of \( X \) the second coordinate is from.

            Similarly, a set \( A \in \mathcal{N} \) is just any any rectangle with one side a measurable set \( F \in \mathcal{M} \) and the other side any set \( E \subset \mathbb{N} \).

            Finally, the measure \(\rho\) just multiplies the measure of the second coordinate set \( F \in \mathcal{M} \) by the number of points in \( E \subset \mathbb{N} \).
          \end{answer}
    \item Describe similarly \( L^+(Y,\rho) \) and \( L_1(Y,\rho) \).
          Rephrase the Fubini and Tonelli theorems in the present context.
          Again, your language should not contain the word ``product,'' but it should explain in more elementary terms the meaning of \( \int_Y f \dd{\rho} \) with \( f \in L^+(Y,\rho) \) or \( f \in L_1(Y,\rho) \), and the meaning of the iterated integrals.
          \begin{answer}
            \( L^+(Y, \rho) \) is the space of non-negative, measurable, two parameter functions \( f: \mathbb{N} \times X \rightarrow \mathbb{R} \).
            This is alternatively the space of sequences \( (f_n) \) of non-negative functions \( f_n: X \rightarrow \mathbb{R} \).
            \( L_1(Y, \rho) \) is then the space of integrable sequences \( (f_n) \) of functions \( f_n: X \rightarrow \mathbb{R} \), i.e., those sequences such that
            \begin{equation}
            \label{eq:5}
            \int |f_n| \dd{\rho} < \infty, \text{ i.e.}, \sum\limits_{n=1}^{\infty} \int |f_n| \dd\mu < \infty
          \end{equation}
          The iterated integrals are
          \begin{equation}
          \label{eq:8}
          \sum\limits_{n=1}^{\infty} \int f_n \dd{\mu}, \quad \int \left(\sum\limits_{n=1}^{\infty} f_n \right) \dd{\mu}.
        \end{equation}

        Tonelli's theorem says that, since every set \( A \in \mathcal{N} \) is \(\sigma\)-finite, that every measurable, non-negative function \( f: \mathbb{N} \times X \rightarrow [0, \infty) \) (sequence of measurable functions) satisfies
        \begin{equation}
        \label{eq:11}
        \int f(n,x) \dd\rho = \sum\limits_{n=1}^{\infty} \int f_n(x) \dd\mu(x) = \int \left( \sum\limits_{n=1}^{\infty} f_n(x) \right) \dd\mu(x).
      \end{equation}
      Fubini's theorem says that, again since every set is \( \sigma \)-finite, for every integrable function \( f: \mathbb{N} \times X \rightarrow \mathbb{R} \), i.e. those sequences such that \cref{eq:5} holds, the sections \( n \) and \( x \) sections (which are a particular function in the sequence and the values for fixed \( x \) of every function in the sequence, respectively), are themselves integrable, and we can interchange the integrals as in \cref{eq:11}.
          \end{answer}
  \end{enumerate}
  \textit{Note:} It is not necessary that you provide proofs.
  A good template is the Background part of the solution to Problem 2.22 (Assignment \#8), where it was explained that a function \( f \) on \( (\mathbb{N}, \nu) \) is a sequence given by \( x_n = f(n) \), the integral of \( f \) is the sum of the series \( \sum\limits_n x_n \), and a function is in \( L_1 \) iff the corresponding series is absolutely convergent.
\end{problem}
\end{document}
